\section{Conclusions}
\label{sec:conclusions}

We introduced an open-weight, open-source framework for translating natural-language strategic requirements into well-formed ATL/ATL$^\ast$ specifications. Three contributions support it: an expert-authored dataset of $1100$ instances spanning simple requirements and five ambiguity families (right and left dislocation, VP ellipsis, Right Node Raising, and quantifier scope ambiguity), balanced over the ATL/ATL$^\ast$ temporal operators and reusable as a benchmark; a hybrid framework running both cloud APIs and local open-weight models with LoRA fine-tuning; and a multi-faceted evaluation protocol coupling normalized exact match with a human-validated LLM-as-a-judge. Our experiments show that in-domain fine-tuning of small, locally deployable open-weight models ($3$--$7$B) translates strategic requirements as accurately as strong few-shot proprietary baselines under the most human-aligned judge; that exact match alone understates quality, since faithful translations often differ syntactically from the reference; and that the automatic judges must themselves be validated, since judge reliability runs opposite to generator capability---the open-weight Llama-3.3-70B is the most human-aligned judge, while the strongest, most recent proprietary models are the least reliable, over-rejecting faithful paraphrases even of their own generations. Finally, the translations plug directly into \textsc{genVITAMIN}, an existing strategic model checker, illustrating a practical deployment path.



\paragraph{Limitations.}
Several limitations qualify these results. The benchmark, although expert-validated, is curated rather than drawn from deployed industrial requirements, and is English-only; transfer to in-the-wild specifications and other languages remains to be demonstrated. We target ATL/ATL$^\ast$, so requirements involving knowledge, beliefs, or the more expressive constructs of Strategy Logic lie outside the present scope. Semantic accuracy is measured by an LLM judge whose agreement with human experts---moderate even for the best-aligned judge---is imperfect, so the absolute numbers are best read as calibrated estimates rather than exact rates; we mitigate this with multiple independent judges and a human-overridden ranking. Our reported confidence intervals, moreover, reflect training-seed variation at a single canonical split; with only $218$ test items, test-set resampling adds uncertainty of order $\pm0.06$ that these intervals do not capture, so they are best read as a lower bound---a stratified $k$-fold evaluation, supported by our framework, would tighten this estimate. Finally, equivalence is assessed by strategic-temporal intent rather than by formal model checking, and the reported cross-environment latencies are only indicative across the API/local boundary.

\paragraph{Future Work.}
Several directions for future work emerge. The main avenue concerns extending the framework to more powerful strategic formalisms, such as Strategy Logic and its fragments~\cite{DBLP:journals/tocl/MogaveroMPV14,DBLP:conf/aaai/CermakLM15,DBLP:conf/ijcai/BelardinelliJKM19} or epistemic variants of ATL~\cite{DBLP:journals/fuin/JamrogaH04,aminof2016prompt,jamroga2017reasoning}, enabling translation about knowledge, beliefs, and information flow in MAS. Moreover, the released dataset opens the door to training domain-specialized models potentially reducing reliance on post-hoc semantic validation.

%Future work will focus on interactive disambiguation, richer strategic formalisms, and deeper coupling between translation models and verification feedback.
